\clearpage
\VerdierPhysicalPageAnchor{0177}%
\begingroup\small\setlength{\parskip}{0pt}\setlength{\emergencystretch}{1.5em}
\setlength{\abovedisplayskip}{4pt plus 2pt minus 2pt}%
\setlength{\belowdisplayskip}{4pt plus 2pt minus 2pt}%

\VerdierSemanticAnchor{ch3d-1-2-1-truncations-cont}
\noindent(resp.
\[
\begin{aligned}
i^q&=\operatorname{identité},&&q<n-1,\\
i^{n-1}&\colon\operatorname{Ker}(d_X^{n-1})\longrightarrow X^{n-1}
  &&\text{(morphisme canonique)},\\
i^q&=0,&&q\geq n
\end{aligned}
\]
).

Remarquons que le morphisme $p$ (resp. $i$) induit un isomorphisme sur les
objets de cohomologie en degré $j\geq n$ (resp. $j<n$) et que
$H^j(X(n,+\infty))=0$ pour $j<n$ (resp. $H^j(X(-\infty,n))=0$ pour
$j\geq n$). Toutes les assertions de la proposition (1.2.1) se démontrent
alors immédiatement en utilisant cette construction. Nous laissons au
lecteur le soin de faire la démonstration en détail.

\VerdierSemanticAnchor{ch3d-1-2-2}
\medskip\noindent\textbf{Définition 1.2.2.}\quad
On appelle \emph{catégorie dérivée de $\mathcal A$} et on note
$D(\mathcal A)$ la catégorie triangulée :
\[
D(\mathcal A)=K(\mathcal A)/\operatorname{Ac}(\mathcal A)
=K(\mathcal A)[\operatorname{Qis}(\mathcal A)^{-1}]
\]
(chap. II, 2.2.6 et 2.2.10).

Il résulte du théorème (2.2.6) que le foncteur cohomologique canonique
$H\colon K(\mathcal A)\longrightarrow\mathcal A$ se factorise d'une manière
unique en un foncteur cohomologique encore noté :

\VerdierSemanticAnchor{ch3d-1-2-2-cohomology}
\begin{equation}
H\colon D(\mathcal A)\longrightarrow\mathcal A
\tag{1.2.2.1}
\end{equation}
et encore appelé \emph{foncteur cohomologique canonique}. Le foncteur
canonique de passage au quotient est noté
$Q\colon K(\mathcal A)\longrightarrow D(\mathcal A)$.

\VerdierSemanticAnchor{ch3d-1-2-3}
\medskip\noindent\textbf{Théorème 1.2.3.}\quad
\textit{a) Les foncteurs naturels qui figurent dans le diagramme :}

\VerdierSemanticAnchor{ch3d-1-2-3-a-diagram}
\begin{center}
  \resizebox{0.80\textwidth}{!}{%
    $\begin{tikzcd}[ampersand replacement=\&,column sep=4.5em,row sep=3.0em]
      \& K^{,+}(\mathcal A)/\operatorname{Ac}(\mathcal A)
        \arrow[dr] \& \\
      K^{,b}(\mathcal A)/\operatorname{Ac}(\mathcal A)
        \arrow[ur] \arrow[dr]
        \& \& D(\mathcal A) \\
      \& K^{,-}(\mathcal A)/\operatorname{Ac}(\mathcal A)
        \arrow[ur] \&
    \end{tikzcd}$%
  }
\end{center}
\endgroup
